\newpage
\begin{center}
	\textbf{\large 5. Инструкция по установке}
\end{center}
% \refstepcounter{chapter}
\addcontentsline{toc}{chapter}{5. Инструкция по установке}

\paragraph{Установка.}
Для сборки и запуска приложения требуется:
\begin{itemize}
	\item компилятор с поддержкой C++17;
	\item CMake версии не ниже 3.15;
	\item OpenCV версии 4.x;
	\item библиотека \texttt{nlohmann\_json}.
\end{itemize}

Исходный код проекта размещён в репозитории GitHub:
\begin{verbatim}
https://github.com/g0lubew-iv/qr-localizer
\end{verbatim}

Для загрузки проекта необходимо клонировать репозиторий и перейти в его директорию:
\begin{verbatim}
git clone https://github.com/g0lubew-iv/qr-localizer
cd qr-localizer
\end{verbatim}

\subsection*{Сборка (Linux/MacOS)}
\paragraph{Шаг 1. Установка зависимостей.}
Установите OpenCV и \texttt{nlohmann\_json} средствами пакетного менеджера вашей системы или с помощью менеджера зависимостей (например, vcpkg).
При использовании Fedora Linux установка может быть выполнена командой:
\begin{verbatim}
sudo dnf install opencv-devel nlohmann-json-devel
\end{verbatim}

\paragraph{Шаг 2. Конфигурация проекта.}
Из корневой директории проекта выполните:
\begin{verbatim}
cmake -S . -B build -DCMAKE_BUILD_TYPE=Release
\end{verbatim}

\paragraph{Шаг 3. Сборка.}
\begin{verbatim}
cmake --build build --config Release
\end{verbatim}

\paragraph{Шаг 4. Установка.}
При необходимости можно установить бинарный файл в системный префикс:
\begin{verbatim}
cmake --install build
\end{verbatim}

\subsection*{Сборка (Windows)}
Сборка под Windows выполняется аналогично с использованием генератора Visual Studio.
Общий подход к конфигурации/сборке/установке для CMake-проектов выполняется по следующей схеме:
\begin{verbatim}
cmake -S . -B build
cmake --build build --config Release
cmake --install build --config Release
\end{verbatim}
При использовании vcpkg необходимо дополнительно указать путь к toolchain-файлу \texttt{vcpkg.cmake} в параметре \verb|-DCMAKE_TOOLCHAIN_FILE=...|.
